\documentclass[11pt,a4paper]{article}
\usepackage[utf8]{inputenc}
\usepackage[margin=2.5cm]{geometry}
\usepackage{amsmath,amssymb,amsfonts}
\usepackage{array}
\usepackage{hyperref}

\hypersetup{
    colorlinks=true,
    linkcolor=blue,
    urlcolor=blue,
    citecolor=blue
}

\setlength{\parindent}{0pt}
\setlength{\parskip}{6pt}

\newcommand{\kotakdefinisi}[2]{%
\noindent\fbox{%
\parbox{\dimexpr\textwidth-2\fboxsep-2\fboxrule\relax}{%
\textbf{#1}\\[4pt]
#2
}%
}\par\vspace{6pt}
}

\begin{document}

% --- Header Catatan Kuliah ---
\begin{center}
    {\LARGE \textbf{CATATAN KULIAH MA4171 TEORI KONTROL LINEAR}}\\[6pt]
    {\Large \textbf{Topik 02: Model Matematika pada Sistem Dinamik}}\\[4pt]
    \textit{Cakupan: Minggu ke-2 sampai Minggu ke-4}\\[2pt]
    \textbf{Referensi Utama:} Katsuhiko Ogata, \textit{Modern Control Engineering} (Bab 2 \& 9.2)
\end{center}
\vspace{-4pt}
\hrule
\vspace{10pt}

\tableofcontents
\vspace{10pt}
\hrule
\vspace{12pt}

\section{Pendahuluan \& Karakteristik Sistem LTI}

Pemodelan matematika merupakan fondasi dasar dalam menganalisis dan merancang sistem kendali. Model matematika merepresentasikan relasi kuantitatif antara masukan (\textit{input}), keluaran (\textit{output}), serta variabel keadaan internal (\textit{state variables}).

\subsection{Karakteristik Sistem Linear Time-Invariant (LTI)}
Kajian utama mata kuliah ini berfokus pada sistem \textbf{Linear Time-Invariant (LTI)} kontinu:
\begin{enumerate}
    \item \textbf{Prinsip Superposisi \& Homogenitas (Linearitas)}:\\
    Jika masukan $u_1(t)$ menghasilkan keluaran $y_1(t)$ dan $u_2(t)$ menghasilkan $y_2(t)$, maka untuk setiap skalar $\alpha, \beta \in \mathbb{R}$:
    \begin{equation*}
        \alpha u_1(t) + \beta u_2(t) \implies \alpha y_1(t) + \beta y_2(t)
    \end{equation*}
    \item \textbf{Time-Invariance (Invarian Waktu)}:\\
    Karakteristik parameter sistem tidak bergantung pada waktu pergeseran. Jika $u(t) \to y(t)$, maka untuk sembarang pergeseran $\tau \ge 0$:
    \begin{equation*}
        u(t - \tau) \implies y(t - \tau)
    \end{equation*}
\end{enumerate}

\subsection{Dua Pendekatan Pemodelan Sistem LTI}
\begin{itemize}
    \item \textbf{Domain Frekuensi (S-Domain / Kontrol Klasik)}: Berbasis \textit{Fungsi Transfer} dan \textit{Diagram Blok}. Pendekatan ini sangat efektif untuk sistem SISO (\textit{Single-Input Single-Output}) dengan asumsi semua kondisi awal bernilai nol.
    \item \textbf{Domain Waktu (State-Space / Kontrol Modern)}: Berbasis sistem persamaan diferensial linear orde satu simultan. Pendekatan ini mampu menangani sistem multivariabel MIMO (\textit{Multi-Input Multi-Output}) dan kondisi awal non-nol.
\end{itemize}

\vspace{6pt}

\section{Fungsi Transfer (\textit{Transfer Function})}

\kotakdefinisi{Definisi Fungsi Transfer:}{
Fungsi transfer sistem LTI didefinisikan sebagai rasio antara Transformasi Laplace dari sinyal keluaran $Y(s)$ terhadap Transformasi Laplace dari sinyal masukan $U(s)$ dengan \textbf{seluruh kondisi awal bernilai nol}:
\begin{equation*}
    G(s) = \frac{Y(s)}{U(s)} = \frac{\mathcal{L}\{y(t)\}}{\mathcal{L}\{u(t)\}}\Bigg|_{\text{seluruh kondisi awal } = 0}
\end{equation*}
}

Jika sistem fisik dimodelkan oleh persamaan diferensial linear orde-$n$:
\begin{equation*}
    a_0 \frac{d^n y}{dt^n} + a_1 \frac{d^{n-1} y}{dt^{n-1}} + \dots + a_{n-1} \frac{dy}{dt} + a_n y(t) = b_0 \frac{d^m u}{dt^m} + b_1 \frac{d^{m-1} u}{dt^{m-1}} + \dots + b_m u(t) \quad (n \ge m)
\end{equation*}
Dengan menerapkan Transformasi Laplace dan menyetel seluruh turunan awal ke nol:
\begin{equation*}
    (a_0 s^n + a_1 s^{n-1} + \dots + a_n) Y(s) = (b_0 s^m + b_1 s^{m-1} + \dots + b_m) U(s)
\end{equation*}
Diperoleh bentuk umum fungsi transfer:
\begin{equation*}
    G(s) = \frac{N(s)}{D(s)} = \frac{b_0 s^m + b_1 s^{m-1} + \dots + b_{m-1} s + b_m}{a_0 s^n + a_1 s^{n-1} + \dots + a_{n-1} s + a_n}
\end{equation*}

\subsection{Karakteristik Penting Fungsi Transfer}
\begin{itemize}
    \item \textbf{Persamaan Karakteristik}: Polinomial penyebut yang disamakan dengan nol, yaitu $D(s) = 0$.
    \item \textbf{Kutub (\textit{Poles})}: Nilai-nilai $s$ yang menyebabkan $D(s) = 0$ ($G(s) \to \infty$). Letak kutub pada bidang kompleks $s$ menentukan kestabilan dan respon alami sistem.
    \item \textbf{Pembuat Nol (\textit{Zeros})}: Nilai-nilai $s$ yang menyebabkan $N(s) = 0$ ($G(s) = 0$).
    \item \textbf{Tanggapan Impuls (\textit{Impulse Response})}: Respon sistem terhadap masukan Dirac delta $\delta(t)$ ($\mathcal{L}\{\delta(t)\} = 1$):
    \begin{equation*}
        y(t) = g(t) = \mathcal{L}^{-1}\{G(s)\}
    \end{equation*}
\end{itemize}

\vspace{6pt}

\section{Aljabar dan Reduksi Diagram Blok}

Diagram blok menyajikan interkoneksi fungsional antar subsistem.

\subsection{Operasi Ekuivalensi Diagram Blok}
\begin{center}
\small
\begin{tabular}{|p{3.5cm}|p{5cm}|p{6cm}|}
\hline
\textbf{Konfigurasi} & \textbf{Deskripsi Hubungan} & \textbf{Fungsi Transfer Ekuivalen} \\ \hline
\textbf{Kaskade (Seri)} & Sinyal mengalir berurutan & $G_{eq}(s) = G_1(s) \cdot G_2(s)$ \\ \hline
\textbf{Paralel} & Sinyal masukan sama dijumlahkan & $G_{eq}(s) = G_1(s) \pm G_2(s)$ \\ \hline
\textbf{Loop Umpan Balik Negatif} & Forward $G(s)$, Feedback $H(s)$ (-) & $T(s) = \dfrac{Y(s)}{R(s)} = \dfrac{G(s)}{1 + G(s)H(s)}$ \\ \hline
\textbf{Loop Umpan Balik Positif} & Forward $G(s)$, Feedback $H(s)$ (+) & $T(s) = \dfrac{Y(s)}{R(s)} = \dfrac{G(s)}{1 - G(s)H(s)}$ \\ \hline
\textbf{Pindah Titik Cabang} & Maju melewati blok $G(s)$ & Jalur cabang dikalikan $\dfrac{1}{G(s)}$ \\ \hline
\textbf{Pindah Titik Jumlah} & Maju melewati blok $G(s)$ & Jalur penjumlahan dikalikan $G(s)$ \\ \hline
\end{tabular}
\end{center}

Untuk sistem dengan umpan balik satuan (\textit{unity feedback}), yaitu $H(s) = 1$:
\begin{equation*}
    T(s) = \frac{G(s)}{1 + G(s)}
\end{equation*}

\vspace{6pt}

\section{Representasi Ruang Keadaan (\textit{State-Space Representation})}

\subsection{Konsep Dasar Ruang Keadaan}
\begin{itemize}
    \item \textbf{Variabel Keadaan (\textit{State Variables})}: Himpunan terkecil dari variabel $\{x_1(t), x_2(t), \dots, x_n(t)\}$ yang bersama dengan masukan $u(t)$ untuk $t \ge t_0$, menentukan dinamika sistem secara utuh untuk setiap $t \ge t_0$.
    \item \textbf{Vektor Keadaan (\textit{State Vector})}: $x(t) = \begin{bmatrix} x_1(t) & x_2(t) & \dots & x_n(t) \end{bmatrix}^T \in \mathbb{R}^n$.
\end{itemize}

\subsection{Bentuk Standar Persamaan Ruang Keadaan LTI}

\kotakdefinisi{Sistem Ruang Keadaan LTI Kontinu:}{
\begin{align*}
    \dot{x}(t) &= A x(t) + B u(t) \quad &\text{(\textbf{Persamaan Keadaan} / \textit{State Equation})} \\
    y(t) &= C x(t) + D u(t) \quad &\text{(\textbf{Persamaan Keluaran} / \textit{Output Equation})}
\end{align*}
}

Dimensi matriks untuk sistem dengan $n$ variabel keadaan, $m$ masukan, dan $p$ keluaran:
\begin{itemize}
    \item $A \in \mathbb{R}^{n \times n}$ : Matriks sistem (\textit{system matrix}).
    \item $B \in \mathbb{R}^{n \times m}$ : Matriks masukan (\textit{input matrix}).
    \item $C \in \mathbb{R}^{p \times n}$ : Matriks keluaran (\textit{output matrix}).
    \item $D \in \mathbb{R}^{p \times m}$ : Matriks transmisi langsung (\textit{direct transmission / feedforward matrix}).
\end{itemize}

\vspace{6pt}

\section{Hubungan Ruang Keadaan dan Fungsi Transfer}

Untuk menurunkan fungsi transfer dari persamaan ruang keadaan, terapkan Transformasi Laplace pada persamaan keadaan dengan asumsi kondisi awal $x(0) = 0$:
\begin{align*}
    s X(s) - x(0) &= A X(s) + B U(s) \\
    (sI - A) X(s) &= B U(s) \\
    X(s) &= (sI - A)^{-1} B U(s)
\end{align*}
Substitusikan ke persamaan Transformasi Laplace dari keluaran $Y(s) = C X(s) + D U(s)$:
\begin{equation*}
    Y(s) = \left[ C(sI - A)^{-1} B + D \right] U(s)
\end{equation*}

\kotakdefinisi{Rumus Konversi Ruang Keadaan ke Fungsi Transfer:}{
\begin{equation*}
    G(s) = \frac{Y(s)}{U(s)} = C(sI - A)^{-1}B + D = \frac{C \operatorname{adj}(sI - A) B}{\det(sI - A)} + D
\end{equation*}
}

\textbf{Keterkaitan Sifat:}\\
Polinomial karakteristik sistem adalah $\det(sI - A) = 0$. Oleh karena itu, nilai-nilai eigen dari matriks $A$ ($\lambda_i$) berkorespondensi langsung dengan kutub-kutub (\textit{poles}) fungsi transfer sistem.

\vspace{6pt}

\section{Bentuk-Bentuk Kanonik Ruang Keadaan}

Diberikan fungsi transfer \textit{strictly proper} ($b_0 = 0$):
\begin{equation*}
    G(s) = \frac{b_1 s^{n-1} + b_2 s^{n-2} + \dots + b_n}{s^n + a_1 s^{n-1} + a_2 s^{n-2} + \dots + a_n}
\end{equation*}

\subsection{1. Bentuk Kanonik Terkontrol (\textit{Controllable Canonical Form})}
Struktur ini sangat penting dalam perancangan kontrol umpan balik keadaan (\textit{state feedback} / penempatan pole):
\begin{equation*}
    A = \begin{bmatrix}
        0 & 1 & 0 & \dots & 0 \\
        0 & 0 & 1 & \dots & 0 \\
        \vdots & \vdots & \vdots & \ddots & \vdots \\
        0 & 0 & 0 & \dots & 1 \\
        -a_n & -a_{n-1} & -a_{n-2} & \dots & -a_1
    \end{bmatrix}, \quad
    B = \begin{bmatrix} 0 \\ 0 \\ \vdots \\ 0 \\ 1 \end{bmatrix}
\end{equation*}
\begin{equation*}
    C = \begin{bmatrix} b_n & b_{n-1} & \dots & b_1 \end{bmatrix}, \quad D = [0]
\end{equation*}

\subsection{2. Bentuk Kanonik Terobservasi (\textit{Observable Canonical Form})}
Bentuk dual dari bentuk terkontrol, sangat fundamental dalam perancangan pengamat keadaan (\textit{state observer}):
\begin{equation*}
    A = \begin{bmatrix}
        0 & 0 & \dots & 0 & -a_n \\
        1 & 0 & \dots & 0 & -a_{n-1} \\
        0 & 1 & \dots & 0 & -a_{n-2} \\
        \vdots & \vdots & \ddots & \vdots & \vdots \\
        0 & 0 & \dots & 1 & -a_1
    \end{bmatrix}, \quad
    B = \begin{bmatrix} b_n \\ b_{n-1} \\ \vdots \\ b_1 \end{bmatrix}
\end{equation*}
\begin{equation*}
    C = \begin{bmatrix} 0 & 0 & \dots & 0 & 1 \end{bmatrix}, \quad D = [0]
\end{equation*}

\subsection{3. Bentuk Kanonik Diagonal / Jordan}
Jika polinomial penyebut memiliki akar-akar real yang berbeda $p_1, p_2, \dots, p_n$:
\begin{equation*}
    G(s) = \frac{c_1}{s - p_1} + \frac{c_2}{s - p_2} + \dots + \frac{c_n}{s - p_n}
\end{equation*}
Maka:
\begin{equation*}
    A = \begin{bmatrix}
        p_1 & 0 & \dots & 0 \\
        0 & p_2 & \dots & 0 \\
        \vdots & \vdots & \ddots & \vdots \\
        0 & 0 & \dots & p_n
    \end{bmatrix}, \quad
    B = \begin{bmatrix} 1 \\ 1 \\ \vdots \\ 1 \end{bmatrix}, \quad
    C = \begin{bmatrix} c_1 & c_2 & \dots & c_n \end{bmatrix}, \quad D = [0]
\end{equation*}

\vspace{6pt}

\section{Transformasi Keserupaan (\textit{Similarity Transformation})}

Representasi ruang keadaan suatu sistem \textbf{tidak bersifat tunggal}. Kita dapat mengubah basis ruang keadaan dengan matriks transformasi non-singular $P \in \mathbb{R}^{n \times n}$ melalui $x(t) = P z(t)$.

\kotakdefinisi{Matriks Baru Hasil Transformasi Keserupaan:}{
\begin{equation*}
    \tilde{A} = P^{-1} A P, \quad \tilde{B} = P^{-1} B, \quad \tilde{C} = C P, \quad \tilde{D} = D
\end{equation*}
}

\textbf{Sifat-Sifat Invarian:}\\
Di bawah transformasi keserupaan:
\begin{enumerate}
    \item Polinomial karakteristik tidak berubah: $\det(sI - \tilde{A}) = \det(sI - A)$.
    \item Nilai-nilai eigen matriks $A$ tidak berubah.
    \item Fungsi transfer $\tilde{G}(s) = G(s)$ identik sama.
\end{enumerate}

\vspace{6pt}

\section{Solusi Domain Waktu \& Matriks Transisi Keadaan}

Untuk sistem keadaan dengan kondisi awal $x(0) = x_0$:
\begin{equation*}
    \dot{x}(t) = A x(t) + B u(t)
\end{equation*}
Solusi analitiknya dalam domain waktu adalah:
\begin{equation*}
    x(t) = e^{At} x(0) + \int_{0}^{t} e^{A(t - \tau)} B u(\tau) \, d\tau
\end{equation*}
\begin{equation*}
    y(t) = C e^{At} x(0) + C \int_{0}^{t} e^{A(t - \tau)} B u(\tau) \, d\tau + D u(t)
\end{equation*}

Di mana \textbf{Matriks Transisi Keadaan} didefinisikan sebagai:
\begin{equation*}
    \Phi(t) = e^{At} = \mathcal{L}^{-1}\{(sI - A)^{-1}\}
\end{equation*}

\textbf{Sifat-Sifat $e^{At}$}:
\begin{itemize}
    \item $e^{A \cdot 0} = I$
    \item $(e^{At})^{-1} = e^{-At}$
    \item $e^{A(t_1 + t_2)} = e^{At_1} e^{At_2} = e^{At_2} e^{At_1}$
    \item $\dfrac{d}{dt} e^{At} = A e^{At} = e^{At} A$
\end{itemize}

\vspace{6pt}

\section{Contoh Soal dan Pembahasan}

\noindent\textbf{Contoh: Menghitung Fungsi Transfer dan Matriks Transisi Keadaan}\\
Diberikan sistem ruang keadaan:
\begin{equation*}
    A = \begin{bmatrix} 0 & 1 \\ -2 & -3 \end{bmatrix}, \quad B = \begin{bmatrix} 0 \\ 1 \end{bmatrix}, \quad C = \begin{bmatrix} 1 & 0 \end{bmatrix}, \quad D = [0]
\end{equation*}

\textit{Penyelesaian:}

Bentuk matriks $(sI - A)$:
\begin{equation*}
    sI - A = \begin{bmatrix} s & -1 \\ 2 & s+3 \end{bmatrix}
\end{equation*}
Determinan $(sI - A)$:
\begin{equation*}
    \det(sI - A) = s(s+3) - (-1)(2) = s^2 + 3s + 2 = (s+1)(s+2)
\end{equation*}
Invers $(sI - A)^{-1}$:
\begin{equation*}
    (sI - A)^{-1} = \frac{1}{(s+1)(s+2)} \begin{bmatrix} s+3 & 1 \\ -2 & s \end{bmatrix}
\end{equation*}

Menghitung fungsi transfer $G(s)$:
\begin{align*}
    G(s) &= C (sI - A)^{-1} B \\
    &= \begin{bmatrix} 1 & 0 \end{bmatrix} \left( \frac{1}{(s+1)(s+2)} \begin{bmatrix} s+3 & 1 \\ -2 & s \end{bmatrix} \right) \begin{bmatrix} 0 \\ 1 \end{bmatrix} \\
    &= \frac{1}{(s+1)(s+2)} \begin{bmatrix} 1 & 0 \end{bmatrix} \begin{bmatrix} 1 \\ s \end{bmatrix} = \frac{1}{s^2 + 3s + 2}
\end{align*}

Menghitung matriks transisi keadaan $\Phi(t) = e^{At} = \mathcal{L}^{-1}\{(sI - A)^{-1}\}$ melalui pecahan parsial:
\begin{align*}
    \mathcal{L}^{-1}\left\{ \frac{s+3}{(s+1)(s+2)} \right\} &= \mathcal{L}^{-1}\left\{ \frac{2}{s+1} - \frac{1}{s+2} \right\} = 2e^{-t} - e^{-2t} \\
    \mathcal{L}^{-1}\left\{ \frac{1}{(s+1)(s+2)} \right\} &= \mathcal{L}^{-1}\left\{ \frac{1}{s+1} - \frac{1}{s+2} \right\} = e^{-t} - e^{-2t} \\
    \mathcal{L}^{-1}\left\{ \frac{-2}{(s+1)(s+2)} \right\} &= \mathcal{L}^{-1}\left\{ \frac{-2}{s+1} + \frac{2}{s+2} \right\} = -2e^{-t} + 2e^{-2t} \\
    \mathcal{L}^{-1}\left\{ \frac{s}{(s+1)(s+2)} \right\} &= \mathcal{L}^{-1}\left\{ \frac{-1}{s+1} + \frac{2}{s+2} \right\} = -e^{-t} + 2e^{-2t}
\end{align*}

Sehingga diperoleh matriks transisi keadaan:
\begin{equation*}
    e^{At} = \begin{bmatrix} 2e^{-t} - e^{-2t} & e^{-t} - e^{-2t} \\ -2e^{-t} + 2e^{-2t} & -e^{-t} + 2e^{-2t} \end{bmatrix}
\end{equation*}

\end{document}
